\documentclass[journal,10pt]{IEEEtran}
\setlength{\textwidth}{182mm}
\setlength{\columnsep}{4.2mm}
\usepackage{amsmath,amssymb,newtxmath}
\usepackage[no-math]{fontspec}
\setmainfont{Times New Roman}
\usepackage{xeCJK}
\setCJKmainfont[AutoFakeBold=2]{SimSun}
\usepackage{graphicx,booktabs,array,hyperref,placeins}
\hypersetup{hidelinks,unicode=true}
\raggedbottom
\setlength{\emergencystretch}{2em}
\renewcommand{\arraystretch}{1.15}
\usepackage{seqsplit}
\begin{document}
\title{Litz-Wire Modeling and Numerical Verification: Comprehensive Research Report · A Mathematical Derivations}
\author{Miao Yuxuan\thanks{Huainan Normal University}}

\maketitle
\noindent 合并稿对应章节 / Chapter in merged manuscript: \href{report_en_ieee.pdf#chapter-A}{A}\par
\FloatBarrier
\section*{A Mathematical Derivations}
\hypertarget{chapter-A}{}
\pdfbookmark[0]{A Mathematical Derivations}{bookmark-chapter-A}
\subsection*{Analytical Normalization, Circuit Constraints, and Element Integrals}
\hypertarget{A.report_derivation.heading}{}
\pdfbookmark[1]{Analytical Normalization, Circuit Constraints, and Element Integrals}{bookmark-A.report_derivation.heading}
Round-wire normalization first requires 2π\(\int\)\(_0\)\(^a\) r Jz(r) dr=It. The identity d[r I\(_1\)(κr)]/dr=κr I\(_0\)(κr) then gives the normalization denominator. Dividing the boundary axial electric field Jz(a)/σ by It yields the internal impedance per unit length. The independent Joule integral is compared with |It|² Re(zint) under the RMS convention. High-precision quadrature error checks only the stated analytical formula; finite ends, external return inductance, temperature changes, and insulation are not included in that analytical accuracy.


Linear-triangle recovery uses the actual physical area aT. With nodal eₖ=v−uₖ, complex element current is σaT\(\sum\)eₖ/3 and Joule power is σaT[\(\sum\)|eₖ|²+|\(\sum\)eₖ|²]/12. The latter follows from the consistent mass matrix with diagonal coefficient two and off-diagonal coefficient one. Original coordinates are scaled by 0.001 m, so area carries the square of that scale; conversion to millimetres for display does not change integration. Material markers distinguish individual strands, return, and air, and the return uses its own v.


Necessity and sufficiency of the common-terminal criterion depend on the same linear system and uniqueness of its solution. Necessity follows by substituting equal currents into the voltage equation. Sufficiency follows because a vector removed by Q can be represented as V1 and then satisfies the augmented system together with total current. Conductor contacts, additional terminal impedances, nonlinear material, or a missing induced term change that system. Reuse of the criterion in a new model therefore begins by reconstructing Z, e, and the constraints, rather than copying an earlier conclusion.


\subsection*{Sharing Derivation}
\hypertarget{A.sharing_derivation.heading}{}
\pdfbookmark[1]{Sharing Derivation}{bookmark-A.sharing_derivation.heading}
Sufficiency of the sharing criterion follows from uniqueness of the port solution. If the projection condition holds, every component of (It/N)Z1+e is equal, so a common voltage exists for which the equal-current vector satisfies both the port equations and the total-current constraint. Necessity follows by substitution of an existing equal-current solution. Numerical sharing is assessed using normalized complex-current deviation, with magnitude and phase reported to retain the information in phasor addition.


\end{document}
